\section{State, rank, and computational cost}
\label{sec:computation}

The closures remove growth with elapsed training time.  They do not remove
every dependence on width or sample count.  Table~\ref{tab:complexity}
separates moving learned state from retained initialized state.

\begin{table}[t]
\centering
\small
\caption{Leading storage and algebraic cost for $L$ hidden layers of width
$n$ and a full batch of $m$ samples.  Constants and first/readout-layer terms
are omitted.  ``Fixed action'' is the multiplication by the retained
$W_{\ell,0}$ and its transpose.}
\label{tab:complexity}
\begin{tabular}{@{}p{0.24\linewidth}p{0.22\linewidth}p{0.22\linewidth}p{0.24\linewidth}@{}}
\toprule
Quantity & Dense flow & Activity clock & Response clock\\
\midrule
Moving internal state
& $(L-1)n^2$
& $2(L-1)nmP$
& $2(L-1)nmP+P^2$\\
Fixed initialized state
& none beyond current weights
& $(L-1)n^2$
& $(L-1)n^2$\\
Learned-correction rank
& at most $n$
& at most $mP$
& at most $m(P+1)$\\
Fixed forward/backward actions, full batch
& included in dense actions
& $O((L-1)n^2m)$
& $O((L-1)n^2m)$\\
Learned factor actions, full batch
& $O((L-1)n^2m)$ total dense work
& $O((L-1)nm^2P)$
& $O((L-1)nm^2P)$ plus solves\\
Small-order algebra
& ---
& $O((L-1)nmP)$ with cumulative triangular sums
& $O(P^3)$ Gram factorization and up to $O((L-1)nmP^2)$ factor solves\\
\bottomrule
\end{tabular}
\end{table}

For Variant A, the triangular dilation sums can be evaluated by cumulative
prefix sums in $O(P)$ per neuron history.  For Variant B, a Cholesky
factorization of the shared Gram matrix costs $O(P^3)$; solving for all factor
coordinates can cost $O((L-1)nmP^2)$ in a direct implementation.  The response
monitor adds directional forward and backward passes.  Conditioning of $G$ can
deteriorate with $P$ even though the prefix proves it remains positive
definite at every fixed order.

The dominant fixed-matrix actions remain quadratic in width.  Consequently,
the present method does not give a subquadratic end-to-end evaluation of a
generic dense initialization.  It can reduce the cost of storing and applying
the \emph{learned correction} when $mP\ll n$, and it can avoid materializing a
dense updated matrix.  A dense implementation that stores only its current
weights may still use less minimal total memory because the closure stores
both $W_0$ and its moving history state.

The theorem turns order into a conservative state--accuracy tradeoff.  Ignoring
constants and thresholds, Variant A requires $P=O(\varepsilon^{-1})$ and hence
$O((L-1)nm\varepsilon^{-1})$ moving history coordinates.  Variant B requires
$P=O(\varepsilon^{-1/2})$, giving
$O((L-1)nm\varepsilon^{-1/2})$ history coordinates plus a
$P^2=O(\varepsilon^{-1})$ shared Gram matrix.  These are upper-bound scalings,
not observed asymptotic rates.  At small practical orders the extra clock and
Gram conditioning can make Variant B less accurate or slower than Variant A.

At large width, each moment column is naturally represented by one value per
neuron.  This is the particle implementation of a population field.  The
current canonical implementation evaluates the learned action through neuron
population averages such as
\begin{equation}
 \frac1nH_{\ell,a,:,k}^\top h_{\ell-1,b},
\end{equation}
and uses the same factors for the transpose action.  Forward and adjoint reuse
of one initialized operator is preserved exactly; redrawing or replacing it
by independent samples would change the training law.

The empirical memory measurements reinforce the distinction.  In the
width-4096 MNIST100 experiment, history factors occupy
$6.25$, $12.50$, and $18.75$ MiB for $P=1,2,3$, compared with $128$ MiB for an
unrestricted learned middle matrix.  The retained $W_0$ adds another
$128$ MiB, so minimal moving-plus-fixed state is slightly larger than the
dense current-state baseline.  The tested closures nevertheless had lower
measured CUDA peaks because their solver workspaces avoided dense evolving
matrix stages; they were slower at the common tolerances.  These are
implementation measurements, not an asymptotic speedup.

